\section{Discussion}\label{sec:discussion}

%Presents a clear statement and insightful interpretation of the results, accurately relates results to the hypothesis, and discusses any limitations or sources of error in the experiment.
%Questions are answered completely and correctly.


The collected data showed that dropping magnets from rest through coils induced an EMF due to a time-varying flux.
The amplitude of each successive coil increased in magnitude as the magnet increased in speed, and the phase switched polarity where the magnets passed through each centre.

By using our derived magnetic dipole moment values from method 1, we calculated a theoretical EMF curve.
This was overlaid on our coil B data as shown in Figures ~\ref{fig:short_coilB} and ~\ref{fig:long_coilB}.
Though the theoretical signal was greater in amplitude compared with our experimental result, the overall shape and
behaviour was consistent with the theory.

\bigskip

The attenuation seen in the experimental result can be attributed to inefficiencies in exchanging the mechanical
energy of the falling magnet with the electrical energy.
Additionally, other magnetic field lines produced from local sources, such as nearby computers and cables, also
affected the EMF
generated by
the
setup.
Further reduction in EMF can be explained by the induced current transferring its energy into Joule heating.

\bigskip

Our analysis on the quadratic interpolation of position showed that the magnets' acceleration was marginally less
than the acceleration due to gravity, $9.81 m/s^2$.
This suggests there is an opposing force to gravity which led to experimental accelerations of $9.41 m/s^2$ and
$9.48 m/s^2$.
We can attribute this to Lenz' Law, where magnetic field generated by the falling magnets resists the direction of
motion.
Instead of using $9.81 m/s^2$ for calculations, the experimental accelerations were used because it
reflected the physical reality more closely.


Due to the geometry of the magnets, the magnetic field lines did not adhere to theoretical $B_z$ when measuring the
field strength near the magnet.
This behaviour changed and $B_z$ began to scale approximately linearly with $z^{-3}$ once the Gauss sensor was far
enough away.
For Figure ~\ref{fig:method1}, the short magnet showed a linear trend for $B$ versus $z^{-3}$, whereas for the majority
of the data points, the
long magnet did not.
The most linear sections of the data were taken in order to find the far field approximation for the magnetic dipole
moment.
This was a problem for the long magnet, as it did not show much linear dependence until further away from the
magnet.
If farther measurements had been taken, the long magnet would have shown a greater linear dependence on $z^{-3}$ and
produced a shallower slope.
This would have been reflected as closer fit to the data for the calculated EMF curve seen in Figure ~\ref{fig:long_coilB}.

\bigskip

Table~\ref{tab:riemann1} and~\ref{tab:riemann2} show the Riemann sum of positive and negative values for
flux across each coil for the short magnet.
This demonstrated that the flux across both phases were approximately equal and opposite for each coil.
The sum total for positive and negative phases was -2.218e-5 Wb and 2.200e-5 Wb, respectively.
This result agrees with the theory since the path of the magnet on either side of each coil is symmetric about the z
axis, and flux is found from Faraday's law, which is dependent on $z$ position.

According to ~\eqref{eq:flux}, there is an upper and lower limit to maximising the magnetic flux.
By taking the numerator, $z^2 - r^2 / 2$, we can determine that if the radius of the coil exceeds $\sqrt{2}z$, then
the flux element starts to become negative.
Therefore, to maximise the net flux as a result of coil radius, we would ideally have a coil with dynamic radius
that would expand or contract to $\sqrt{2}z$.
The lower limit is the radius of the magnet itself.
However, the radial element increases at the same rate as $B_z$ decreases, thus cancelling the scaling effect.
This also explains why the calculated flux is symmetrical for both positive and negative phases despite the
acceleration of the
magnet.

\section{Conclusion}\label{sec:conclusion}

%Presents a clear, concise, focused, insightful and correct connection between the results and the investigated
%concepts.
%Explains why the findings are important, what they show, how they contribute to the broader understanding of the topic and recommendations for future research or improvements to
%experimental design.
%The section brings closure to the report.

It was discovered that a moving magnet produces a time-varying flux and thus an EMF through
multi-turn
coils.
The resulting induced current also form its own magnetic field which opposed the direction of the field lines of the
magnet, agreeing with Lenz' Law.
Method 1 yielded magnetic dipoles of \qty{0.56}{\ampere\meter^2} and \qty{2.05}{\ampere\meter^2} for the short and
long magnets,
respectively.
Though these values resulted in overshooting the calculated amplitude of the EMF versus the experimental signal, the
theory showed an agreement with the overall shape of the response curve.

Deviations from the calculated signal were attributed to other local magnetic field sources, geometric limitations, and
Joule heating.
For better results, the experiment should be performed in a container which shields itself from other magnetic fields.